\chapter{Evaluating and selecting a detector}
\label{ch:eval}

A detector's numbers mean nothing unless the data split is honest and the
metrics are the ones an operator would recognise. This chapter describes the
split in \code{pyplant/ml/dataset.py}, the metrics in
\code{pyplant/ml/evaluate.py}, and the rule that picks the detector to
deploy.

\section{Split by scenario, never by sample}

Windows overlap, and neighbouring windows from one run are nearly identical. A
random sample split puts window 41 in training and window 42 in test; the
model scores brilliantly on data it has effectively already seen, and the
number is not slightly optimistic, it is meaningless. Whole runs are split
instead.

\begin{figure}[htbp]
\centering
\scalebox{0.8}{%
\begin{tikzpicture}[node distance=5mm and 14mm]
  \node[box,minimum width=40mm] (all) {all scenarios};
  \node[good,below left=8mm and 4mm of all,minimum width=34mm] (nor) {normal runs\\(shuffled, seed)};
  \node[bad,below right=8mm and 4mm of all,minimum width=34mm] (fau) {faulty runs};
  \node[box,below=10mm of nor,xshift=-34mm] (tr) {train\\60\%};
  \node[box,right=of tr] (ca) {calibrate\\20\%};
  \node[box,right=of ca] (tn) {test normal\\20\%};
  \node[box,right=of tn] (tf) {test faulty\\100\%};
  \draw[flow] (all) -- (nor);
  \draw[flow] (all) -- (fau);
  \draw[flow] (nor) -- (tr);
  \draw[flow] (nor) -- (ca);
  \draw[flow] (nor) -- (tn);
  \draw[flow] (fau) -- (tf);
  \node[below=1mm of tr,font=\footnotesize] {fit model};
  \node[below=1mm of ca,font=\footnotesize] {set threshold};
  \node[below=1mm of tn,font=\footnotesize] {false alarms};
  \node[below=1mm of tf,font=\footnotesize] {detection, delay, isolation};
\end{tikzpicture}}
\caption{\code{split_scenarios()}. Normal runs are divided three ways; faulty
runs are used only for testing.}
\label{fig:split}
\end{figure}

Each pool has a single job. Using one pool for fitting, thresholding and false
alarm measurement flatters the threshold. At least three normal runs are
required; the Quick preset makes six and the Standard preset ten. Because each
grid entry has its own seed, magnitude and onset, held-out runs differ from
training runs in the ways that matter operationally.

\section{Persistence}

A single flagged window is not a detection. A detection is declared at the
first window of the first run of $p = 3$ consecutive flagged windows
(\code{PERSISTENCE}). With a 10\,s stride this costs at most 20\,s of delay.

\begin{figure}[htbp]
\centering
\scalebox{0.8}{%
\begin{tikzpicture}[x=6mm,y=6mm]
  \foreach \i/\f in {0/0,1/1,2/0,3/0,4/1,5/1,6/0,7/0,8/1,9/1,10/1,11/1,12/1}{
    \ifnum\f=1 \fill[badred!60] (\i,0) rectangle (\i+0.9,0.9);
    \else \fill[panelgrey] (\i,0) rectangle (\i+0.9,0.9); \draw[linegrey] (\i,0) rectangle (\i+0.9,0.9); \fi
  }
  \draw[okgreen!70!black,very thick] (7.95,-0.1) rectangle (10.95,1.0);
  \node[font=\footnotesize,below] at (9.45,-0.15) {detection declared here (3 in a row)};
  \node[font=\footnotesize,above] at (1.45,1.0) {flicker: ignored};
  \node[font=\footnotesize,above] at (4.95,1.0) {2 in a row: ignored};
  \node[font=\footnotesize,left] at (0,0.45) {window flags};
\end{tikzpicture}}
\caption{The persistence rule. Isolated flags never count, either as
detections or as false alarms.}
\label{fig:persist}
\end{figure}

A detector that flickers above threshold one window in twenty has a
fine-looking detection rate and is unusable. The same rule counts false
alarms, so both numbers describe the same detector.

\section{The metrics}

\subsection{False alarms per day}
Measured on the held-out normal runs:
\begin{equation}
  \text{FAR} = \frac{\text{number of persistent excursions}}
                    {\text{normal time covered (s)}} \times 86\,400 .
\end{equation}
Operators experience alarms per shift, not per sample. The report also gives
the raw fraction of normal windows above threshold, which should be close to
$1-q$.

\subsection{Detection rate and delay}
For each faulty test run, the detection delay is the time from the declared
onset to the first persistent detection after it:
\begin{equation}
  \tau_d = t_\text{detect} - t_\text{onset} .
\end{equation}
Detection rate is the fraction of faulty runs detected; the report gives the
median delay overall and per fault mode. Incipient faults (long ramps) have
naturally longer delays.

\subsection{Isolation accuracy}
Once a fault has fully developed (ramp finished plus
\code{DEVELOPED_AFTER} $= 300$\,s), the most frequent top channel over the
remaining windows is compared with the channel physics names
(Chapter~\ref{ch:ml}). Judging earlier is wrong: fouling ramping over twenty
minutes is still a transient five minutes in, and the transient implicates
the controlled level rather than the downstream one.

If a run ends before the fault has developed, isolation is judged on the last
quarter of the post-onset data and the run is marked as not fully developed.

\subsection{Consistency}
For each mode, the fraction of runs whose top channel is that mode's most
common top channel. It needs no ground truth, so it is the only isolation
number available for sensor faults.

\begin{warn}[Why sensor faults are not scored for isolation]
In a closed loop, a level transmitter reading 25\,mm high is driven back to
setpoint by the controller: the reading looks perfect and the true level sits
25\,mm low. The signature appears everywhere \emph{except} the faulty
instrument, so ``the top channel should be \tg{LT101}'' is exactly backwards.
Scoring against that expectation would be scoring against a mistake.
\end{warn}

\section{Selecting the detector}

Each chosen method is fitted, calibrated and evaluated on the same split. The
one deployed is
\begin{equation}
  \arg\max_m \big(\text{detection rate}_m,\; -\text{FAR}_m\big)
\end{equation}
in lexicographic order: more detections wins, a tie goes to fewer false
alarms, and remaining ties go to the simpler model because the methods are
tried simplest first (z-score, PCA, isolation forest, autoencoder). A learned
detector that cannot beat PCA is not deployed just because it is newer.

\begin{figure}[htbp]
\centering
\scalebox{0.8}{%
\begin{tikzpicture}[node distance=7mm]
  \node[term] (s) {methods ticked on \ui{Model}};
  \node[box,below=of s] (f) {for each: fit, calibrate, evaluate};
  \node[decision,below=of f] (d) {highest\\detection rate?};
  \node[decision,below=10mm of d] (t) {tie: fewest\\false alarms?};
  \node[decision,below=10mm of t] (u) {still tied:\\simplest method};
  \node[good,below=of u] (o) {\code{artifacts/detector.joblib}\\+ \code{manifest.json}, \code{metrics.json}};
  \draw[flow] (s) -- (f);
  \draw[flow] (f) -- (d);
  \draw[flow] (d) -- (t);
  \draw[flow] (t) -- (u);
  \draw[flow] (u) -- (o);
\end{tikzpicture}}
\caption{Model selection in \code{train_detector()}.}
\label{fig:select}
\end{figure}

\section{Reference results}

With the full dataset (58 one-hour runs: \code{python -m pyplant.ml.train
--duration 3600}) the project's reference numbers are:

\begin{center}
\begin{tabular}{@{}lcccc@{}}
\toprule
& Detection & Median delay & False alarms & Isolation \\
\midrule
PCA (MSPC baseline) & 83.3\% & 256\,s & 0 / day & 60\% \\
Autoencoder & \textbf{100\%} & \textbf{70\,s} & \textbf{0 / day} & \textbf{73\%} \\
\bottomrule
\end{tabular}
\end{center}

PCA detects only a third of the sensor bias and drift runs; the autoencoder
detects all of them. The physics channels do the isolation: leak runs point at
\tg{FLOW_IMBALANCE} 12 out of 12 times, pump wear at \tg{HEAD_RESIDUAL} 6/6
and stiction at \tg{OP_MINUS_ZT} 6/6. The raw window flag rate on normal runs
is 0.14\% and there are no persistent excursions. The z-score and isolation
forest detectors were added later; run them on your own dataset to place them
in this table. Smaller presets give noisier numbers.

\section{Reading the Model tab}

\begin{enumerate}
  \item The comparison table has one row per method: detection rate, median
        delay, false alarms per day, flagged normal windows and isolation.
        The deployed method is marked \emph{(selected)}.
  \item The per-mode table shows, for the selected method, the number of runs
        of each fault mode, the fraction detected, the median delay, the most
        common top channel and the channel physics expects.
  \item A low detection rate on one mode with a long delay usually means the
        fault was still ramping when the run ended; use longer runs.
  \item A non-zero false alarm rate means the threshold is too low for this
        data: raise the quantile (e.g.\ 0.999) or generate more normal runs.
\end{enumerate}
